\chapter{Les LLM au service de l'administration des réseaux}
\label{chap:partie-2}

La partie précédente a montré qu'un réseau peut être automatisé sans être autonome. Reste à déterminer quel type d'intelligence artificielle peut combler cet écart. Les grands modèles de langage (LLM, \emph{Large Language Models}) présentent à cet égard un atout certain : ils traitent du texte, et l'administration des réseaux en est largement composée. Configurations, journaux d'événements, documentations constructeur et tickets d'incident relèvent tous, sous des formes diverses, du langage.

Les études examinées conduisent à une réponse nuancée à notre problématique. Un LLM rend des services réels pour lire, expliquer et aider au diagnostic. En revanche, sa fiabilité reste insuffisante pour agir seul sur un réseau, car il ne peut être sûr des conséquences de ses propositions. Nous présentons d'abord le fonctionnement d'un LLM et les moyens de l'adapter au domaine des réseaux, puis les effets du dialogue en langage naturel. Nous examinons ensuite les usages étudiés, avant de les confronter aux résultats d'une évaluation en conditions réelles.

\section{Fonctionnement des LLM et adaptation au domaine des réseaux}

\paragraph{Principe.}
Un LLM est un modèle d'apprentissage profond fondé sur l'architecture Transformer \cite{vaswani2017}. Grâce à un mécanisme d'attention, il évalue l'importance de chaque mot d'un texte au regard des autres. L'entraînement consiste à prédire le jeton suivant d'une séquence, c'est-à-dire un mot ou une fraction de mot, tâche répétée sur des volumes de texte considérables. À partir d'une certaine taille, le modèle réalise des tâches pour lesquelles il n'a pas été entraîné explicitement, à partir de quelques exemples fournis dans la requête. Pour un administrateur, un même outil peut ainsi expliquer une configuration, résumer des journaux ou proposer une commande.

\paragraph{Connaissances et lacunes.}
Une partie de la documentation publique sur les réseaux (RFC, manuels, forums) figure dans les données d'entraînement. Un modèle généraliste connaît donc les protocoles courants. Ses lacunes demeurent toutefois importantes. Il ne dispose d'aucune information sur l'infrastructure qui lui est confiée, qu'il s'agisse de la topologie, des configurations ou des incidents en cours. Ses connaissances sont figées à la date de son entraînement et ne sont fiables que pour les sujets documentés publiquement. Lors de l'expérience ShowNet présentée en section~\ref{sec:evaluations}, le modèle a régulièrement produit des configurations erronées pour des équipements dont la documentation n'est pas publique. Il a même confondu deux systèmes d'exploitation Cisco aux syntaxes différentes \cite{nakamura2026}.

\paragraph{Langage naturel et langage réseau.}
Les données d'un réseau diffèrent du texte ordinaire. Une liste de contrôle d'accès (ACL), une commande ou un modèle YANG (\emph{Yet Another Next Generation}) obéissent à des règles formelles et associent symboles, adresses et contraintes numériques \cite{huang2024}. Une phrase telle que « restreindre l'accès au serveur 192.168.1.5 depuis toutes les adresses externes » doit être traduite en une règle d'ACL exacte, sans approximation possible. Certaines données, comme une topologie ou des mesures de télémétrie, ne se présentent pas sous forme textuelle et doivent être converties. Dans l'expérience ShowNet, la topologie du cœur de réseau a ainsi été décrite en notation Mermaid, directement dans les instructions du modèle \cite{nakamura2026}.

La question du volume se pose également. En 2023, la fenêtre de contexte des modèles, c'est-à-dire la quantité de texte qu'ils peuvent traiter simultanément, suffisait à peine à décrire une topologie de quelques dizaines de nœuds \cite{mani2023}. Elle s'est fortement accrue depuis, mais un réseau d'opérateur, avec ses configurations et ses journaux, demeure bien trop volumineux pour être fourni dans son intégralité.

\paragraph{Techniques d'adaptation.}
Réentraîner un modèle depuis le début dépasse les moyens de la plupart des organisations. On combine donc plusieurs techniques plus légères \cite{liu2025,hong2025}. La figure~\ref{fig:adaptation} situe ces techniques dans la chaîne qui relie la demande de l'opérateur au réseau.

\begin{figure}[htbp]
\centering
\begin{tikzpicture}[node distance=0.5cm and 0.5cm]
  \node[human, text width=1.7cm] (op) {Opérateur};
  \node[llm, right=of op, text width=2.6cm] (req) {Requête structurée\\[-1pt]{\scriptsize\textit{prompt engineering}}};
  \node[llm, right=of req, text width=2.3cm, minimum height=1.5cm, line width=1pt] (llm) {\textbf{LLM}\\[-1pt]{\scriptsize éventuellement \textit{fine-tuné}}};
  \node[box, right=of llm, text width=2.4cm] (prop) {Proposition\\[-1pt]{\scriptsize réponse, configuration ou code}};
  \node[human, right=of prop, text width=2.1cm] (val) {Validation\\[-1pt]{\scriptsize vérificateur ou humain}};
  \node[tool, above=0.8cm of llm, text width=3.9cm] (rag) {Base documentaire (RAG)\\[-1pt]{\scriptsize documentation, normes, configurations}};
  \node[tool, below=0.8cm of llm, text width=3.9cm] (tools) {Outils via MCP\\[-1pt]{\scriptsize ligne de commande, supervision, tickets}};
  \node[net, text width=2.1cm] (net) at (val |- tools) {Réseau};
  \draw[arr] (op) -- (req);
  \draw[arr] (req) -- (llm);
  \draw[arr] (llm) -- (prop);
  \draw[arr] (prop) -- (val);
  \draw[arr] (rag) -- node[lab, right]{contexte} (llm);
  \draw[arr2] (llm) -- node[lab, right]{appels} (tools);
  \draw[arr2] (tools) -- node[lab, above]{lecture de l'état} (net);
  \draw[arr] (val) -- node[lab, right]{application} (net);
\end{tikzpicture}
\caption{Place des techniques d'adaptation entre la demande de l'opérateur et le réseau. Le prompt engineering agit sur la requête, le RAG ajoute du contexte, le fine-tuning modifie le modèle et les outils donnent accès à l'état réel du réseau. Dans tous les cas, le LLM formule une proposition, qui doit être validée avant d'être appliquée.}
\label{fig:adaptation}
\end{figure}

La technique la plus simple est le prompt engineering, qui consiste à soigner la requête : rôle attendu, contexte, exemples et format de la réponse. Les effets peuvent être considérables. Pour l'analyse de journaux, des stratégies de requête élaborées ont amélioré les performances jusqu'à 380,7~\% par rapport à une requête naïve \cite{liu2024}. Le RAG (\emph{Retrieval-Augmented Generation}) ajoute une recherche documentaire avant chaque réponse \cite{lewis2020}. Le système extrait de la documentation constructeur, des normes ou des configurations en place les passages pertinents, qu'il transmet au modèle avec la question. La réponse s'appuie alors sur des informations à jour et propres au réseau.

Le fine-tuning modifie le modèle lui-même, en poursuivant son entraînement sur des données du domaine. Des méthodes légères en réduisent aujourd'hui le coût, mais son efficacité dépend de la qualité des données. Un modèle spécialisé devient en outre obsolète s'il n'est pas mis à jour avec l'évolution des technologies \cite{huang2024}. L'appel d'outils permet enfin au modèle d'interroger lui-même des équipements ou des interfaces de programmation (API, \emph{Application Programming Interface}). Pour Huang et al., un LLM destiné aux réseaux doit combiner ces deux compétences : comprendre le langage et savoir utiliser des outils \cite{huang2024}. Le \emph{Model Context Protocol} (MCP) tend à s'imposer comme standard de cet accès. Le chatbot de l'expérience ShowNet s'en servait pour interroger les routeurs et le système de tickets \cite{nakamura2026}. Cette capacité prépare les agents présentés dans la partie suivante.

Une dernière approche consiste à demander au LLM d'écrire le programme qui traitera les données, plutôt que de lui confier les données elles-mêmes. Les données restent chez l'opérateur, ce qui limite les risques pour la confidentialité, les adresses IP étant des données personnelles. Le coût ne dépend plus de la taille du réseau, et l'opérateur peut relire le code avant de l'exécuter. Sur l'analyse du trafic, cette méthode atteint en moyenne 68~\% de réponses justes sur quatre modèles, et 88~\% avec GPT-4, contre 23~\% lorsque le graphe est fourni directement au modèle \cite{mani2023}.

\section{Compréhension du langage naturel et interaction humain-réseau}

\paragraph{Un besoin ancien.}
Un administrateur réseau alterne en permanence entre plusieurs langages : lignes de commande propres à chaque constructeur, modèles YANG, outils de supervision, scripts internes. Sans approche unifiée, l'apprentissage est long et les erreurs sont fréquentes \cite{mani2023}. Parler au réseau en langage courant n'est pas une idée nouvelle. Dès 2021, le système Lumi permettait à un opérateur de saisir une intention du type « inspecter le trafic des résidences étudiantes ». Il en extrayait les éléments utiles par apprentissage automatique, les traduisait dans un langage intermédiaire, sollicitait la confirmation de l'opérateur pour vérifier sa compréhension, puis déployait la configuration. Dans l'étude utilisateur, 88,5~\% des participants déclaraient préférer Lumi, seul ou en complément des commandes \cite{jacobs2021}. Ces systèmes reposaient toutefois sur des modèles entraînés pour une tâche précise. Ils généralisaient mal, coûtaient cher à entraîner et s'intégraient difficilement à d'autres outils \cite{huang2024}. Les LLM permettent de lever ces trois obstacles grâce à un modèle généraliste unique.

\paragraph{Une interaction dans les deux sens.}
L'administrateur peut demander au modèle de traduire une consigne en commandes ou en configuration. Il peut aussi lui demander l'inverse : rendre intelligible ce que le réseau produit, qu'il s'agisse d'expliquer une configuration, de résumer une série d'alertes ou de reconstituer un incident. Les usages observés lors de l'expérience ShowNet, détaillés en section~\ref{sec:evaluations}, montrent que ces deux directions ne sont pas utilisées de la même façon selon le profil de l'utilisateur. Pour les premiers, le chatbot jouait le rôle d'un tuteur. Pour les seconds, celui d'un assistant. Les explications produites sont jugées satisfaisantes. Six praticiens, dont l'expérience dépasse dix ans, ont attribué en moyenne 4,42 sur 5 à la lisibilité et à l'utilité des analyses de journaux produites par LogPrompt \cite{liu2024}.

\paragraph{Lien avec l'Intent-Based Networking.}
On retrouve ici l'IBN (\emph{Intent-Based Networking}) présenté dans la première partie. La RFC 9315 définit une intention comme un objectif formulé de manière déclarative, sans préciser la manière de l'atteindre \cite{clemm2022}. Le langage naturel constitue le moyen le plus direct d'exprimer un tel objectif. Hong et al. citent d'ailleurs l'interprétation des intentions parmi les principaux atouts des LLM pour la gestion des réseaux \cite{hong2025}. La figure~\ref{fig:ibn} indique où intervient le modèle dans ce processus, ainsi que les limites de son rôle.

\begin{figure}[htbp]
\centering
\begin{tikzpicture}[node distance=0.4cm]
  \node[human, text width=2.2cm] (i1) {Intention en langage naturel};
  \node[llm, right=of i1, text width=2.2cm] (i2) {Levée d'ambiguïté};
  \node[llm, right=of i2, text width=2.2cm] (i3) {Politique formelle};
  \node[llm, right=of i3, text width=2.2cm] (i4) {Configuration};
  \node[tool, right=of i4, text width=2.2cm] (i5) {Vérification};
  \node[ex, below=0.15cm of i1, text width=2.4cm, anchor=north] (e1) {\og interdire SSH depuis le Wi-Fi invités\fg{}};
  \node[ex, below=0.15cm of i2, text width=2.4cm, anchor=north] (e2) {Quelles interfaces ? Avant ou après les règles existantes ?};
  \node[ex, below=0.15cm of i3, text width=2.4cm, anchor=north] (e3) {refuser TCP/22 des invités vers les équipements};
  \node[ex, below=0.15cm of i4, text width=2.4cm, anchor=north] (e4) {ACL appliquée aux bonnes interfaces};
  \node[ex, below=0.15cm of i5, text width=2.4cm, anchor=north] (e5) {vérificateur ou tests sur le réseau};
  \node[human, above=1.2cm of i2, text width=2.2cm] (op) {Opérateur};
  \draw[arr] (i1) -- (i2);
  \draw[arr] (i2) -- (i3);
  \draw[arr] (i3) -- (i4);
  \draw[arr] (i4) -- (i5);
  \draw[arr] (op.west) -| (i1.north);
  \draw[arr2] (op) -- node[lab, right]{questions\\réponses} (i2);
  \draw[arr, dashed] (i5.north) |- node[lab, above, pos=0.75]{intention respectée ?} (op.east);
  \node[legbox, fill=blue!12, draw=blue!60!black, anchor=west] (l1) at ($(i1.west |- e2.south)+(0,-0.45)$) {};
  \node[right=2pt of l1, font=\scriptsize] (t1) {étape confiée au LLM};
  \node[legbox, fill=green!10, draw=green!45!black, right=0.6cm of t1] (l2) {};
  \node[right=2pt of l2, font=\scriptsize] (t2) {outil classique};
  \node[legbox, fill=orange!15, draw=orange!70!black, right=0.6cm of t2] (l3) {};
  \node[right=2pt of l3, font=\scriptsize] {humain};
\end{tikzpicture}
\caption{Rôle du LLM dans une chaîne d'Intent-Based Networking, illustré par l'accès SSH. Le modèle traduit l'intention, mais doit d'abord lever les ambiguïtés en interrogeant l'opérateur \cite{mondal2025}. La vérification finale reste confiée à des outils classiques, conformément à la RFC 9315 \cite{clemm2022}.}
\label{fig:ibn}
\end{figure}

\paragraph{Le problème de l'ambiguïté.}
Le langage naturel présente un défaut majeur : il est ambigu. Demander de « prioriser les visioconférences » ne précise ni les applications concernées, ni le degré de priorité à leur accorder, ni la partie du réseau visée. Mondal et al. montrent que la difficulté persiste même lorsque le modèle traduit correctement chaque règle. Une nouvelle règle d'ACL ou de routage recouvre souvent des règles existantes, et leur ordre de priorité modifie le comportement du réseau d'une manière que le modèle ne peut pas deviner seul. Dans un grand réseau de cloud, les auteurs ont relevé des ACL présentant plusieurs centaines de recouvrements. Leur prototype, Clarify, ajoute au LLM un module de désambiguïsation qui interroge l'opérateur avant de générer la configuration, puis de la vérifier \cite{mondal2025}. Cette boucle de questions est représentée sur la figure~\ref{fig:ibn}. Interpréter une intention relève donc davantage du dialogue que de la traduction automatique.

\section{Cas d'usage, du plus simple au plus exigeant}

Les états de l'art \cite{liu2025,hong2025} recensent de nombreux usages. Pour les organiser, nous les regroupons selon l'exigence qu'ils imposent au modèle : comprendre un contenu existant, produire un contenu nouveau ou raisonner en plusieurs étapes (figure~\ref{fig:gradation}). À mesure que l'on progresse, une erreur du modèle a des conséquences plus lourdes, et la validation devient plus exigeante.

\begin{figure}[htbp]
\centering
\begin{tikzpicture}[col/.style={draw, rounded corners=3pt, align=center, font=\footnotesize,
                    text width=4.2cm, minimum height=4.1cm, inner sep=5pt}]
  \node[col, fill=blue!6, draw=blue!40!black] (c1) {%
    {\normalsize\textbf{Comprendre}}\\[1pt]\textit{le LLM lit}\\[6pt]
    Analyse de journaux\\ Documentation et questions\\ Explication de configuration\\[6pt]
    \textcolor{black!60}{\rule{3.4cm}{0.3pt}}\\[3pt]
    \textbf{Validation :} l'opérateur relit};
  \node[col, fill=blue!14, draw=blue!50!black, right=0.45cm of c1] (c2) {%
    {\normalsize\textbf{Produire}}\\[1pt]\textit{le LLM écrit}\\[6pt]
    Génération de configuration\\ Traduction d'intention\\ \phantom{x}\\[6pt]
    \textcolor{black!60}{\rule{3.4cm}{0.3pt}}\\[3pt]
    \textbf{Validation :} vérificateur puis humain};
  \node[col, fill=blue!24, draw=blue!60!black, right=0.45cm of c2] (c3) {%
    {\normalsize\textbf{Raisonner}}\\[1pt]\textit{le LLM enchaîne des étapes}\\[6pt]
    Diagnostic\\ Planification\\ \phantom{x}\\[6pt]
    \textcolor{black!60}{\rule{3.4cm}{0.3pt}}\\[3pt]
    \textbf{Validation :} ingénieur d'astreinte};
  \draw[-{Stealth[length=3mm]}, line width=1.6pt, blue!50!black]
    ([yshift=-0.35cm]c1.south west) -- node[below=2pt, font=\scriptsize, text=black]{exigence demandée au modèle et coût d'une erreur}
    ([yshift=-0.35cm]c3.south east);
\end{tikzpicture}
\caption{Classement des cas d'usage selon l'exigence qu'ils imposent au modèle. Plus la tâche est exigeante, plus la validation humaine ou outillée doit être poussée.}
\label{fig:gradation}
\end{figure}

\paragraph{Comprendre.}
L'analyse de journaux est l'usage le plus immédiat. Les équipements produisent un volume important de journaux, aux formats variés et souvent mal structurés, que les outils classiques traitent au moyen de règles écrites à la main ou de modèles entraînés pour un format précis. LogPrompt, évalué sur des journaux de systèmes logiciels, ne nécessite aucun entraînement sur le domaine. Il obtient néanmoins jusqu'à 55,9~\% de performances supérieures à celles de méthodes entraînées sur des milliers de journaux \cite{liu2024}. La principale limite tient au volume : il y a bien trop de journaux pour les confier intégralement au modèle, et un tri préalable reste indispensable. L'explication de configurations relève de la même logique, comme le montre l'expérience ShowNet, où le modèle a identifié les paramètres manquants d'une configuration BGP (\emph{Border Gateway Protocol}) \cite{nakamura2026}.

\paragraph{Documenter et répondre aux questions.}
Répondre à une question sur une technologie suppose de s'appuyer sur la documentation disponible. C'est le rôle du RAG présenté plus haut. Dans l'expérience ShowNet, le RAG seul a été jugé utile dans 61,4~\% des cas, et il atteint 79,1~\% lorsqu'il est combiné à l'accès à la ligne de commande \cite{nakamura2026}.

\paragraph{Produire.}
La génération de configurations est l'usage le plus étudié, et les premiers résultats appellent la prudence. Mondal et al. ont constaté que GPT-4 livré à lui-même produit des brouillons prometteurs, mais comportant de nombreuses erreurs de topologie, de syntaxe et de sémantique. Le résultat change lorsque le modèle est placé dans une boucle de vérification qui lui renvoie des erreurs précises (figure~\ref{fig:boucle}). Il corrige alors la plupart de ses erreurs seul, et l'humain n'intervient qu'en dernier recours \cite{mondal2023}.

\begin{figure}[htbp]
\centering
\begin{tikzpicture}[node distance=0.45cm]
  \node[human, text width=2.1cm] (ex) {Exigence\\[-1pt]{\scriptsize ex. traduire Cisco vers Juniper}};
  \node[llm, right=of ex, text width=1.4cm, minimum height=1.2cm, line width=1pt] (m) {\textbf{LLM}};
  \node[box, right=of m, text width=2.2cm] (c) {Configuration candidate};
  \node[tool, right=of c, text width=2.4cm] (v) {Vérificateurs\\[-1pt]{\scriptsize syntaxe et sémantique}};
  \node[box, right=1.2cm of v, text width=2.2cm, fill=green!22, draw=green!45!black] (ok) {Configuration vérifiée};
  \node[human, above=0.8cm of m, text width=1.6cm] (h) {Humain};
  \draw[arr] (ex) -- (m);
  \draw[arr] (m) -- (c);
  \draw[arr] (c) -- (v);
  \draw[arr] (v) -- node[lab, above]{conforme} (ok);
  \draw[arr] (v.south) -- ++(0,-0.7) -| node[lab, below, pos=0.25]{erreur localisée renvoyée au modèle} (m.south);
  \draw[arr, dashed] (h) -- node[lab, right]{en dernier recours} (m);
\end{tikzpicture}
\caption{Génération de configuration avec boucle de vérification, d'après Mondal et al. \cite{mondal2023}. Dans la traduction de Cisco vers Juniper, les auteurs comptent environ dix retours automatiques pour une intervention humaine. Pour une politique de non-transit, ce rapport passe à six.}
\label{fig:boucle}
\end{figure}

La traduction d'intentions repose sur le même principe, auquel s'ajoute l'étape de levée d'ambiguïté (figure~\ref{fig:ibn}). L'intention « interdire l'accès SSH aux équipements depuis le Wi-Fi invités » doit donner lieu à une liste de contrôle d'accès insérée au bon endroit parmi les règles existantes \cite{mondal2025}. Le système IBN doit ensuite vérifier que l'intention est effectivement respectée \cite{clemm2022}.

\paragraph{Raisonner.}
Le diagnostic est plus exigeant. Il suppose de formuler des hypothèses, de choisir les vérifications à effectuer, d'interpréter les résultats et, parfois, de remettre en cause ses premières conclusions. RCACopilot, développé chez Microsoft, collecte automatiquement les informations de diagnostic d'un incident, puis demande au LLM de prédire la catégorie de la cause et de justifier sa réponse. Évalué sur un an d'incidents réels, il atteint un Micro-F1 de 0,766 \cite{chen2024}. Le LLM assiste ainsi l'ingénieur d'astreinte sans jamais intervenir lui-même sur le réseau. Le raisonnement en plusieurs étapes reste cependant fragile : dans l'expérience ShowNet, le modèle a échoué à une requête combinant deux conditions, qu'il a tenté de traiter en une seule commande \cite{nakamura2026}.

La planification va plus loin. Pour la mise à jour logicielle d'un parc de routeurs, le modèle peut découper le travail en étapes, depuis la vérification de l'état initial jusqu'au contrôle final, en prévoyant un plan de retour arrière. Le système ChatNet applique cette démarche à un cas de planification de capacité, mais les auteurs reconnaissent que le modèle n'égale pas encore un expert \cite{huang2024}.

\section{Ce que montrent les évaluations en conditions réelles}
\label{sec:evaluations}

Décrire ce que les LLM savent faire ne suffit pas à répondre à notre problématique. Il faut aussi observer ce qu'ils font réellement. L'étude la plus détaillée dont nous disposons est celle menée sur ShowNet, le réseau de démonstration du salon Interop Tokyo. Le réseau comptait environ 700 équipements, dont des routeurs et des commutateurs de huit constructeurs. Pendant deux semaines, 105 ingénieurs ont utilisé un chatbot fondé sur GPT-4.1. Celui-ci pouvait consulter la documentation du réseau, interroger les équipements en ligne de commande via MCP et lire les tickets. Par précaution, la modification des configurations avait été désactivée : le modèle pouvait lire l'état du réseau, mais pas le modifier \cite{nakamura2026}.

\begin{figure}[htbp]
\centering
\begin{minipage}[b]{0.52\textwidth}
\centering
\begin{tikzpicture}
\begin{axis}[barres, width=4.3cm, height=3.6cm, title={Par type de tâche},
  yticklabels={Supervision, Explications techniques, Analyse de configuration, Dépannage, Gestion de projet}]
  \addplot[fill=blue!35, draw=blue!60!black] coordinates {(77.0,1) (70.2,2) (68.9,3) (58.7,4) (50.0,5)};
\end{axis}
\end{tikzpicture}
\end{minipage}\hfill
\begin{minipage}[b]{0.44\textwidth}
\centering
\begin{tikzpicture}
\begin{axis}[barres, width=3.6cm, height=2.2cm, title={Par outil utilisé},
  yticklabels={RAG seul, Ligne de commande seule, RAG et ligne de commande}]
  \addplot[fill=green!40, draw=green!45!black] coordinates {(61.4,1) (73.4,2) (79.1,3)};
\end{axis}
\end{tikzpicture}
\end{minipage}
\caption{Part des réponses jugées utiles par les participants de l'expérience ShowNet, parmi les réponses évaluées \cite{nakamura2026}. Le taux global est de 68,1~\%.}
\label{fig:shownet}
\end{figure}

Sur les réponses évaluées, 68,1~\% ont été jugées utiles. La figure~\ref{fig:shownet} montre toutefois que ce taux varie fortement : 77~\% pour la supervision, contre 58,7~\% pour le dépannage. Sur 1\,427 commandes envoyées, 85,1~\% ont été exécutées sans erreur de syntaxe, sur des équipements dont la documentation est accessible en ligne. Les usages diffèrent selon le profil. Les ingénieurs juniors s'en sont surtout servis pour se faire expliquer des technologies (44~\% de leurs échanges) et pour analyser des configurations (35~\%). Les membres expérimentés de l'équipe d'exploitation l'utilisaient davantage pour la supervision et le dépannage \cite{nakamura2026}. Les spécialistes envoyés par les constructeurs n'ont jugé utiles que 43,8~\% des réponses, ce qui s'explique en partie par des demandes qui excédaient les capacités de l'outil.

Le modèle a aussi produit des erreurs graves. Il a inventé, dans une de ses réponses, une liste de tickets inexistants \cite{nakamura2026}. Ces erreurs confirment que toute sortie doit être vérifiée avant d'être appliquée.

\paragraph{Portée des résultats.}
Ces résultats doivent être lus avec prudence. Il s'agit d'une seule étude, menée sur un réseau d'exposition construit pour une démonstration, et non sur un réseau de production. Les taux reposent sur des évaluations déclarées par les participants, et non sur une mesure objective de la justesse des réponses. Nous n'avons pas trouvé, à ce stade, d'autre étude de terrain comparable pour les confronter.

\paragraph{Déploiements en production.}
RCACopilot assiste le diagnostic d'incidents chez Microsoft, sans intervenir sur le réseau \cite{chen2024}. Ce cas illustre le niveau d'autonomie le plus répandu aujourd'hui : le LLM analyse et conseille, et l'humain décide. Les systèmes plus autonomes, qui agissent eux-mêmes sur l'infrastructure, sont présentés dans la partie suivante.

\paragraph{Bilan.}
En production, le LLM demeure un assistant étroitement encadré. Ses droits sont limités, il reçoit des données déjà collectées, ou ses propositions sont soumises à validation. L'écart entre tâches simples et tâches complexes est net : la supervision et l'explication fonctionnent mieux que le diagnostic ou les opérations en plusieurs étapes. Enfin, l'efficacité de l'outil dépend largement de son utilisateur. Un opérateur qui connaît les limites du modèle en tire davantage de profit que celui qui en attend des solutions miracles.

Si les LLM suscitent un tel intérêt pour l'administration des réseaux, c'est d'abord parce qu'ils traitent des données textuelles très variées sans qu'il soit nécessaire d'entraîner un modèle pour chaque tâche, contrairement aux approches d'apprentissage automatique présentées dans la partie précédente. C'est ensuite parce qu'ils permettent de dialoguer avec le réseau en langage courant, ce qui rapproche l'administration de l'expression d'objectifs. Les évaluations rappellent toutefois qu'un LLM seul se limite à proposer, qu'il s'agisse d'une configuration, d'une hypothèse ou d'un plan. Il ne perçoit le réseau qu'à travers ce qu'on lui communique, ne peut agir dessus qu'à l'aide d'outils et n'offre aucune garantie quant à la justesse de ses productions. Les deux premières limites conduisent naturellement aux agents et aux systèmes multi-agents, traités dans la partie suivante. La troisième concerne la sécurité et le contrôle humain, qui feront l'objet d'une partie dédiée.
